\section{结论}
本文提出 PPK-FIP，将 TMP 酯化的指纹选择性估计转化为有限宽度先验，并在共享 Ping Pong 结构中显式加入伪持续性与批内一阶失活。相对于重新估计的初始简化模型和初始全模型，PPK-FIP 的合并留一时间点预测误差分别降低 \LotoAllRedOrig\,\% 和 \LotoAllRedFull\,\%。消融结果将主要预测改善归于 Bornadel 数据的伪持续性修正，以及稀疏曲线中的指纹先验约束。已知真值模拟表明，显式伪持续性能够减小固定指纹选择性造成的估计偏差。

PPK-FIP-GP 通过高斯过程残差校正进一步降低了训练窗口内的预测误差，但前半段训练的外推中，PPK-FIP 与初始简化模型的合并误差接近。动力学模型的内插精度、外推能力与参数机理解释应分别评价。对于生物润滑剂过程设计，指纹先验和显式机理扩展可作为有限数据条件下的建模选择，参数应用仍需受数据覆盖范围、传质条件和外部验证结果约束。
